%%%%%%%%%%%%%%%%%
% Configuration %
%%%%%%%%%%%%%%%%%

% Spec: 2-5 pages (excluding references). Do not change font, margins, or spacing.
\documentclass[12pt, letterpaper]{article}
\usepackage{xurl}
\usepackage[super,comma,sort&compress]{natbib}
\usepackage{amsmath}
\usepackage{graphicx}
\usepackage{listings}
\lstset{basicstyle=\ttfamily\small, columns=fullflexible}
\usepackage{lipsum}

% Any configuration that should be done before the end of the preamble:
\usepackage{hyperref}
\hypersetup{colorlinks=true, urlcolor=blue, linkcolor=blue, citecolor=blue}

\renewcommand{\refname}{References Cited}

\title{ECS 240 Project Proposal}
\author{Rajiv Choudhury, Jacob Hamzawi, Rishika Nigam, Gurpreet Singh}

\begin{document}

\maketitle

%%%%%%%%%%%
% Article %
%%%%%%%%%%%

\section{Introduction and Motivation}

\subsection{Problem Domain and Motivation}
% TODO: Application domain; what problem(s) we solve and why it's interesting.
% At least 3 citations, at least 2 of them research papers.


%TODO: tidy up
Functional programming allows for greater modularity.\cite{hughes_why_1989}
Side effects cause difficulty in reasoning about code; unclear when calling a function which components are modified.\cite{cousot_purity_2005}
Yet effects are often absolutely necessary as with file I/O.
Desire to break free from assignment and mutation \cite{backus_can_1978}

\subsection{Proposed Language}
% TODO: Own toy language, or a subset of an existing one? What it contains/does,
% or which keywords/features we focus on. Include example syntax.

We propose a toy programming language similar to traditional systems programming languages like C and Rust with the distinction that every function is explicitly defined as "pure" or "impure", where a "pure" function may only modify state initialized in or moved into its scope.
Our language will implement at least the minimum set of elements that allow us to illustrate this system, including variables, functions, references, and an 8-bit unsigned integer primitive as the sole type.
%TODO: fix
We may support other features such as structs, strings, or simple file I/O depending on our timeline (see \S\S~2~and~4).

Rust itself originally sought such a design, defining all functions as "pure" by default and employing the \texttt{io} keyword to explicitly designate a function with I/O side effects.\cite{hoare2013pure}
The language designers eventually opted to remove \texttt{io} keyword and the pure/impure separation.
%TODO: make sure we define side effects before using it here
We hope to discover while designing and analyzing our own language what explicit declaration of side effects could enable for language users and whether the Rust developers' decision was ultimately justified.

\begin{verbatim}
fn pure_function(x, y) {
    x + y
}

// global variable
var z;

// You must specify the data out of scope whose
// state may change
fn impure(z) impure_function(x, y) {
    z = x + y;
}

// If variables have fields, those fields
// must be declared 
fn impure(x.value, y.value) impure_function_struct(x, y) {
    // mutates x
    x.value = 1;
    y.value = 1;
}

// You can look at the program see a full list
// of all the things that could change state
fn impure main() {
    var x = 1;
    var y = 2;
    x = pure_function(x, y);
	  // mutated: z
	  impure_function(x, y);
}
\end{verbatim}

\subsection{Proposed Tool}
% TODO: Basic interpreter + one of: compiler, static analysis, dynamic analysis,
% type system. What are the inputs/outputs? How does a user interact with it?

compiler, where the input is a string in our source language and the output is a string in Rust.
The user interacts with it from the command line
As a stretch goal (see \S2), an analysis tool to glean information from impure declarations may be in scope.

\subsection{Scope and Related Work}
% TODO: Is this already solved/built? What existing tools solve similar problems?

Haskell has the Monad \cite{peyton_jones_imperative_1993}
Functional programming allows for greater modularity.\cite{hughes_why_1989}
And they also often discourage side effects by the nature of their design.
\texttt{cargo-scan} is a static analysis tool for Rust that detects side effects from Rust source code.\cite{zoghbi_auditing_2026}

\section{Goals}
% TODO: High level; details go in later sections.

% I kinda discussed above; we want to learn how to track scope, detect impure functions, enumerate all possible mutations from a function, and in general, discover whether a pure/impure designation would be helpful in a systems programming language.

\paragraph{Minimum goal.}
% Must have working by end of quarter.

Parser, interpreter, and working compiler.
Interpreter and compiler must detect whether a function is pure or impure and error if a function is impure without declaring so along with the outside state it modifies.

\paragraph{Target goal.}
% Plan to have working by end of quarter.

The above, along with a static analysis tool that enumerates all the possible state that is modified given a function (including the functions it itself calls).
The tool will at least be a CLI, but could be extended to a VSCode extension as well.

\paragraph{Reach goal.}
% Extra features if time allows; nothing on the critical path.

%TODO:
The above, as well as other Rust features that might make the division harder and mitigation: system state (file I/O), structs, etc.?
Specifically for file I/O, we could explore options to avoid Rust's issue of \texttt{io} keyword explosion (similar to Haskell's Monad).

\section{Syntax and Semantics}
% TODO: Example syntax, draft formal grammar, and at most a paragraph sketching
% the semantics (set of (INPUT, prog, OUTPUT) triples).

\begin{lstlisting}
example program goes here
\end{lstlisting}

% TODO:
\begin{align*}
  e &::= n \mid x \mid e_1 + e_2 \\
  c &::= x := e \mid c_1; c_2
\end{align*}

\section{Implementation Plan}

\subsection{Tool and Language Selection}
% TODO: What tools and languages will we use to build the project?

Parser in Rust; meta language is Rust; target language for compiler is Rust.

\subsection{Incorporating Existing Tools}
% TODO: From scratch, or using existing software/code bases? Justify each one.

From scratch.

\subsection{Architecture}
% TODO: Components and how they interact. Consider a diagram (can be tentative).

Core language definition; parser that outputs core; interpreter reads core and executes; compiler reads core and outputs IR, in IR we do checks on purity, then IR to target Rust

% \begin{figure}[h]
%   \centering
%   \includegraphics[width=\linewidth]{architecture.pdf}
%   \caption{Tool architecture.}
%   \label{fig:architecture}
% \end{figure}

\subsection{Development Effort}
% TODO: Functions/components/modules to build (refer to architecture).
% What is explicitly out of scope?

\subsection{Challenges}
% TODO: What could go wrong? What's the fallback if a step gets difficult?

% TODO:
Tracking scope is hard; tool is too much work...
Fallback..?
What if a variable will \emph{never} change state? if we have if statements...should we detect that? need it be declared?

\section{Collaboration Plan}
% TODO: How are components divided among members? Meeting plan (suggested
% 2x/week) and communication platform.

% TODO:
\paragraph{Division of Work.} Work divided as follows

\bigskip

\begin{tabular}{ll}
  \textbf{Component} & \textbf{Group Member(s)} \\
  \hline
  Core & Jacob Hamzawi \\
  Parser & ? \\
  Interpreter & ? \\
  Compiler & ? \\
  Analysis tool & ? \\
\end{tabular}

\paragraph{Communication.}Slack and GitHub; Meetings after class on Monday and over Zoom on Thursday

\paragraph{Timeline.}
% Expected weekly timeline of which components are developed when.

The below timeline assumes a 4-week time allotment to complete the project.
It may be contracted or expanded depending on our eventual project due date.

% TODO:
\begin{tabular}{ll}
  \textbf{Week} & \textbf{Components} \\
  \hline
  Week 1 & Parser \\
  Week 2 & Interpreter + Compiler \\
  Week 3 & Analysis tool + exploration \\
  Week 4 & Presentation \\
\end{tabular}

\paragraph{Commitment statement.}

% TODO: uncomment when each person agrees
% \begin{itemize}
%   \item I commit to work and contribute equally to this project. --- \emph{Rajiv Choudhury}
%   \item I commit to work and contribute equally to this project. --- \emph{Jacob Hamzawi}
%   \item I commit to work and contribute equally to this project. --- \emph{Rishika Nigam}
%   \item I commit to work and contribute equally to this project. --- \emph{Gurpreet Singh}
% \end{itemize}

\section{Conclusion and Questions for Feedback}
% TODO: 1-2 paragraphs. What happens if successful, why it's exciting/important.
% What are we uncertain about / want feedback on?

%TODO:
Re-state goal kinda: does explicit knowledge of impure functions change programming paradigm? Or something like that.

Questions: are variables enough?
Should we focus instead on system state?
That's more difficult but perhaps more compelling.

%%%%%%%%%%%%%%
% References %
%%%%%%%%%%%%%%

\bibliographystyle{unsrtnat}
\bibliography{references}

% Material beyond the 5-page limit goes here (may or may not be graded):
% \appendix
% \section{...}

\end{document}
